\chapter{The Spiral Geometry of Reality}

In RCM, spirals are the fundamental geometrical structures that organize coherent motion and memory. This chapter explores why spirals minimize action, how they couple linear and angular momentum, and how nested spiral hierarchies emerge.



%---------------------------------------------------------------
\section{Why Spirals? Mathematical Motivation}

\paragraph{Spiral efficiency principle.}  
Whenever a body must \emph{translate} through space while it \emph{rotates} about an axis, the most energy‐efficient (least‐action) trajectory is a helix.  By wrapping rotation around translation, a spiral lets one unit of motion satisfy two needs at once:
\begin{itemize}
  \item Translational kinetic energy drives forward progress.
  \item The same motion supplies the centrifugal radius required by the spin, so no extra curvature energy (or radiation) is wasted.
\end{itemize}
The Resonance Continuum Model (RCM) elevates this to a law: coherent structures—from photons to hurricanes—appear as spirals precisely because spirals minimise integrated kinetic energy while storing a memory of their path.

\subsubsection{Least‐Action Derivation}

Consider a point mass \(m\) moving at constant speed \(v\) along a helix of radius \(r\), carrying axial spin \(L_z=I\,\omega\) (moment of inertia \(I\), angular speed \(\omega\)).  The kinetic action over a time interval \(\Delta t\) is
\[
  S(r)
  = \int_{0}^{\Delta t}
      \Bigl[\tfrac12\,m v^{2} \;+\; \tfrac12\,I\,\omega^{2}\Bigr]\,dt.
\]
They obey the constraint
\[
  L_z = I\,\omega = m\,v\,r,
  \tag{2.1}
\]
so eliminating \(\omega\) gives the kinetic energy per unit time
\[
  T(r)
  = \tfrac12\,m\,v^{2}
    + \frac{L_z^{2}}{2\,m\,v\,r^{2}},
  \quad
  S(r) = T(r)\,\Delta t.
\]
Setting \(\partial S/\partial r=0\) yields the optimal radius
\[
  \boxed{\,r = \frac{L_z}{m\,v}\,}.
  \tag{2.2}
\]

\subsubsection{Photon Helicity as a Helix}

For a circularly‐polarised photon \(L_z=\hbar\) and \(v=c\), using the effective mass
\[
  m = \frac{h}{\lambda\,c}
  \quad(E = h\nu)
\]
in (2.2) gives
\[
  r = \frac{\hbar}{(h/\lambda c)\,c}
     = \frac{\lambda}{2\pi}.
\]
Thus a photon’s wavefront can be pictured as a helix of radius \(\lambda/2\pi\): one full optical period advances one azimuthal turn—the most efficient way for light to carry its spin.

\subsubsection{Coupling Linear and Angular Momentum}

Along a helix,
\[
  p = m\,v,
  \qquad
  L = m\,v\,r,
\]
and since \(r=L/(mv)\) is fixed by (2.2), \(p\) and \(L\) are \emph{geometrically locked}: changing one forces the other to adjust, or energy is wasted in curvature.

\bigskip
\begin{table}[H]
  \small
  \centering
  \caption{Spiral systems across scales}
  \label{tab:spiral-systems}
  \begin{tabular}{@{}p{0.20\textwidth} p{0.25\textwidth} p{0.25\textwidth} p{0.30\textwidth}@{}}
    \toprule
    \textbf{System (scale)} & \textbf{Tangential momentum \(p\)} & \textbf{Axial spin \(L\)} & \textbf{Spiral consequence} \\
    \midrule
    Cyclotron electron     & \(m\,v\)                     & \(m\,v\,r\)           & Radius set by \(B\)-field \\
    Magnetic skyrmion core & \(\hbar\,k\)                 & \(n\hbar\)            & Quantised helix \\
    Hurricane eyewall      & \(\rho\,v\)                  & \(\rho\,v\,r\)        & Pitch fixed by Coriolis \\
    \bottomrule
  \end{tabular}
\end{table}

\bigskip

\subsubsection{Nested Spiral Hierarchies}

Coarse‐graining a tight helix over many turns produces a looser spiral at a larger scale.  Repeating this zoom‐out builds a \emph{fractal ladder}:
\[
  \text{DNA double helix}
  \;\to\;
  \text{protein super‐helices}
  \;\to\;
  \text{neuronal scroll waves}
  \;\to\;
  \text{atmospheric vortices}
  \;\to\;
  \text{spiral galaxies}.
\]
Each tier obeys \(r=L/(mv)\) with effective collective parameters, furnishing RCM’s multiscale memory backbone.

\subsubsection{Empirical Manifestations of Spiral Efficiency}

\begin{itemize}
  \item \textbf{Optical Bessel beams} propagate as self‐healing helices, minimising diffraction loss.
  \item \textbf{Theta–gamma spiral lattices} in the hippocampus pack temporal codes at minimal metabolic cost.
  \item \textbf{Black‐hole ring-downs} settle into logarithmic spirals in the complex‐frequency plane, shedding excess curvature energy until the helix implied by the Kerr metric is reached.
\end{itemize}

\subsubsection{Key Points}

\begin{enumerate}
  \item \textbf{Efficiency law:} For fixed spin and speed, a helix of radius \(r=L/(mv)\) minimises action.
  \item \textbf{Photon benchmark:} A photon’s radius \(r=\lambda/2\pi\) is the smallest that accommodates \(\hbar\) of spin without extra energy.
  \item \textbf{Momentum locking:} Linear and angular momentum are inseparable on a helix; decoupling them wastes energy.
  \item \textbf{Multiscale economy:} Nested spirals propagate this efficiency upward, enabling coherent, low-loss structures from molecules to galaxies.
\end{enumerate}

With spiral efficiency established, Section 2.2 will classify specific spiral families—logarithmic, Archimedean, power-law, composite—and compute their minimum resonance frequencies within the RCM framework.


%---------------------------------------------------------------
\subsection{Spiral Sub-classes and Their Minimum Resonance Frequencies}

Section 2.1 explained *why* helices appear; here we identify *which* spiral a resonant bundle adopts and derive the *lowest frequency* each class can sustain while remaining coherent.

\subsubsection{Pitch, Curvature, and the Resonance Cut-off}

For any planar spiral \(r(\theta)\) (polar radius vs.\ azimuthal angle) define
\begin{align}
  p(\theta)
  &= \frac{1}{r}\,\frac{dr}{d\theta},
  &
  \kappa(\theta)
  &= \frac{\bigl[r^{2} + 2\bigl(dr/d\theta\bigr)^{2} - r\,d^{2}r/d\theta^{2}\bigr]^{1/2}}
           {\bigl[r^{2} + \bigl(dr/d\theta\bigr)^{2}\bigr]^{3/2}}.
\end{align}
RCM’s amplitude–frequency law (\(A\omega=\mathrm{const}\)) implies that a spiral decoheres if its intrinsic rotation rate drops below the curvature-induced floor:
\begin{equation}
  \boxed{\;\omega_{\min}(\theta)=v\,\kappa(\theta)\;},
  \tag{2.3}
\end{equation}
with \(v\) the bundle’s center-of-mass speed.  For each spiral family we (i) write \(r(\theta)\), (ii) compute \(p(\theta)\) and \(\kappa(\theta)\), and (iii) find the minimum of (2.3).

\subsubsection*{Logarithmic spiral}

\begin{equation}
  r(\vartheta) = a\,e^{b\vartheta} \qquad (\text{constant pitch } p=b).
\end{equation}

\begin{itemize}[leftmargin=1.5em]
  \item Pitch: \(p(\vartheta)=b\).
  \item Curvature: \(\displaystyle \kappa(\vartheta)
           = \frac{b}{\sqrt{1+b^{2}}}\,e^{-b\vartheta}
           \;\xrightarrow[\vartheta\to\infty]{}\;0.\)
\end{itemize}

\begin{equation}
  \boxed{\;\omega_{\min}=0 \quad (\text{formally})\;}
\end{equation}

\emph{Interpretation:} a logarithmic spiral can extend arbitrarily far; coherence is
noise-limited, not geometry-limited.

\subsubsection*{Archimedean spiral}

\begin{equation}
  r(\vartheta) = a + b\,\vartheta \qquad (\text{linear growth}).
\end{equation}

\begin{itemize}[leftmargin=1.5em]
  \item Pitch: \(\displaystyle p(\vartheta)=\frac{b}{\,a+b\vartheta\,}.\)
  \item Curvature: \(\displaystyle
    \kappa(\vartheta)
    = \frac{a + 2b\vartheta}
           {\bigl[(a+b\vartheta)^{2}+b^{2}\bigr]^{3/2}}.\)
\end{itemize}

For a finite disc of radius \(R\), the lowest sustainable rotation scales as
\begin{equation}
  \boxed{\;
    \omega_{\min}
    \;=\; C\,\frac{v\,b^{3}}{(R - b)^{2}}
    \;\propto\; R^{-2}
  \;}
\end{equation}
with \(C>0\) a shape-dependent constant.


\subsubsection*{Power-law spiral}

\begin{equation}
  r(\vartheta) = a\,\vartheta^{\,n}, \qquad n>0.
\end{equation}

\begin{itemize}[leftmargin=1.5em]
  \item Pitch: \(p(\vartheta)=\dfrac{d\ln r}{d\vartheta}=\dfrac{n}{\vartheta}\).
  \item Curvature (scaling): \(\kappa(\vartheta)\propto \vartheta^{-(n+1)}\).
\end{itemize}

Given an outer radius \(R\), the outer turn is
\begin{equation}
  \vartheta_{\max}=\Bigl(\tfrac{R}{a}\Bigr)^{\!1/n}.
\end{equation}
If \(v\) is the characteristic speed along the spiral, the low–frequency cut-off scales as
\begin{equation}
  \boxed{\;
    \omega_{\min}
    \;=\; C_{n}\,\frac{v\,a^{\,1+1/n}}{R^{\,1+1/n}}
    \;},
\end{equation}
where \(C_n>0\) is a shape-dependent constant. Note: as \(n\to 0^{+}\), \(r(\vartheta)\to a\) (a circle); at \(n=1\) one recovers the Archimedean spiral \(r=a\vartheta\). A logarithmic spiral belongs to the distinct family \(r=a\,e^{b\vartheta}\).

\subsubsection*{Composite spirals}\chapter{The Spiral Geometry of Reality}

In \RCM, spirals are the fundamental geometrical structures that organize coherent motion and memory. This chapter explores why spirals minimize action, how they couple linear and angular momentum, and how nested spiral hierarchies emerge.

%---------------------------------------------------------------
\section*{Why spirals? Mathematical motivation}

\paragraph{Spiral efficiency principle.}
Whenever a body must \emph{translate} through space while it \emph{rotates} about an axis, the most energy-efficient (least-action) trajectory is a helix. By wrapping rotation around translation, a spiral lets one unit of motion satisfy two needs at once:
\begin{itemize}[leftmargin=1.5em]
  \item Translational kinetic energy drives forward progress.
  \item The same motion supplies the centrifugal radius required by the spin, so no extra curvature energy (or radiation) is wasted.
\end{itemize}
The Resonance Continuum Model elevates this to a law: coherent structures—from photons to hurricanes—appear as spirals precisely because spirals minimize integrated kinetic energy while storing a memory of their path.

\subsubsection*{Least-action derivation}

Consider a point mass \(m\) moving at constant speed \(v\) along a helix of radius \(r\), carrying axial spin \(L_z=I\,\omega\) (moment of inertia \(I\), angular speed \(\omega\)). The kinetic action over a time interval \(\Delta t\) is
\[
  S(r)
  = \int_{0}^{\Delta t}
      \Bigl[\tfrac12\,m v^{2} \;+\; \tfrac12\,I\,\omega^{2}\Bigr]\,dt.
\]
They obey the constraint
\begin{equation}
  L_z \;=\; I\,\omega \;=\; m\,v\,r,
  \label{eq:helix-constraint}
\end{equation}
so eliminating \(\omega\) gives the kinetic energy per unit time
\[
  T(r)
  = \tfrac12\,m\,v^{2}
    + \frac{L_z^{2}}{2\,m\,v\,r^{2}},
  \qquad
  S(r) \;=\; T(r)\,\Delta t.
\]
Setting \(\partial S/\partial r=0\) yields the optimal radius
\begin{equation}
  \boxed{\,r \;=\; \frac{L_z}{m\,v}\,}.
  \label{eq:opt-radius}
\end{equation}

\subsubsection*{Photon helicity as a helix}

For a circularly polarized photon \(L_z=\hbar\) and \(v=c\), using the effective mass
\[
  m \;=\; \frac{h}{\lambda\,c}
  \qquad (E = h\nu)
\]
in \eqref{eq:opt-radius} gives
\[
  r \;=\; \frac{\hbar}{(h/\lambda c)\,c}
     \;=\; \frac{\lambda}{2\pi}.
\]
Thus a photon’s wavefront can be pictured as a helix of radius \(\lambda/2\pi\): one full optical period advances one azimuthal turn—the most efficient way for light to carry its spin.

\subsubsection*{Coupling linear and angular momentum}

Along a helix,
\[
  p \;=\; m\,v,
  \qquad
  L \;=\; m\,v\,r,
\]
and since \(r=L/(mv)\) is fixed by \eqref{eq:opt-radius}, \(p\) and \(L\) are \emph{geometrically locked}: changing one forces the other to adjust, or energy is wasted in curvature.

\begin{table}[H]
  \small
  \centering
  \caption{Spiral systems across scales.}
  \label{tab:spiral-systems}
  \renewcommand{\arraystretch}{1.15}
  \begin{tabularx}{\linewidth}{@{}p{0.20\linewidth} p{0.25\linewidth} p{0.25\linewidth} X@{}}
    \toprule
    \textbf{System (scale)} & \textbf{Tangential momentum \(p\)} & \textbf{Axial spin \(L\)} & \textbf{Spiral consequence} \\
    \midrule
    Cyclotron electron     & \(m\,v\)                     & \(m\,v\,r\)           & Radius set by \(B\)-field \\
    Magnetic skyrmion core & \(\hbar\,k\)                 & \(n\hbar\)            & Quantized helix \\
    Hurricane eyewall      & \(\rho\,v\)                  & \(\rho\,v\,r\)        & Pitch fixed by Coriolis \\
    \bottomrule
  \end{tabularx}
\end{table}

\subsubsection*{Nested spiral hierarchies}

Coarse-graining a tight helix over many turns produces a looser spiral at a larger scale. Repeating this zoom-out builds a \emph{fractal ladder}:
\[
  \text{DNA double helix}
  \;\to\;
  \text{protein super-helices}
  \;\to\;
  \text{neuronal scroll waves}
  \;\to\;
  \text{atmospheric vortices}
  \;\to\;
  \text{spiral galaxies}.
\]
Each tier obeys \(r=L/(mv)\) with effective collective parameters, furnishing \RCM’s multiscale memory backbone.

\subsubsection*{Empirical manifestations of spiral efficiency}

\begin{itemize}[leftmargin=1.5em]
  \item \textbf{Optical Bessel beams} propagate as self-healing helices, minimizing diffraction loss.
  \item \textbf{Theta–gamma spiral lattices} in the hippocampus pack temporal codes at minimal metabolic cost.
  \item \textbf{Black-hole ring-downs} settle into logarithmic spirals in the complex-frequency plane, shedding excess curvature energy until the helix implied by the Kerr metric is reached.
\end{itemize}

\subsubsection*{Key points}

\begin{enumerate}[leftmargin=1.5em,label=\arabic*.]
  \item \textbf{Efficiency law:} For fixed spin and speed, a helix of radius \(r=L/(mv)\) minimizes action.
  \item \textbf{Photon benchmark:} A photon’s radius \(r=\lambda/2\pi\) is the smallest that accommodates \(\hbar\) of spin without extra energy.
  \item \textbf{Momentum locking:} Linear and angular momentum are inseparable on a helix; decoupling them wastes energy.
  \item \textbf{Multiscale economy:} Nested spirals propagate this efficiency upward, enabling coherent, low-loss structures from molecules to galaxies.
\end{enumerate}

With spiral efficiency established, the next section classifies specific spiral families—logarithmic, Archimedean, power-law, composite—and computes their minimum resonance frequencies within the \RCM\ framework.

%---------------------------------------------------------------
\section*{Spiral sub-classes and their minimum resonance frequencies}

\paragraph{Pitch, curvature, and the resonance cut-off.}
For any planar spiral \(r(\vartheta)\) (polar radius vs.\ azimuthal angle) define
\begin{align}
  p(\vartheta)
  &= \frac{1}{r}\,\frac{dr}{d\vartheta},
  &
  \kappa(\vartheta)
  &= \frac{\bigl[r^{2} + 2\bigl(dr/d\vartheta\bigr)^{2} - r\,d^{2}r/d\vartheta^{2}\bigr]^{1/2}}
           {\bigl[r^{2} + \bigl(dr/d\vartheta\bigr)^{2}\bigr]^{3/2}}.
\end{align}
\RCM’s amplitude–frequency law (\(A\omega=\mathrm{const}\)) implies that a spiral decoheres if its intrinsic rotation rate drops below the curvature-induced floor:
\begin{equation}
  \boxed{\;\omega_{\min}(\vartheta)=v\,\kappa(\vartheta)\;},
  \label{eq:omega-min}
\end{equation}
with \(v\) the bundle’s center-of-mass speed. For each spiral family we (i) write \(r(\vartheta)\), (ii) compute \(p(\vartheta)\) and \(\kappa(\vartheta)\), and (iii) find the minimum of \eqref{eq:omega-min}.

\subsubsection*{Logarithmic spiral}

\begin{equation}
  r(\vartheta) = a\,e^{b\vartheta} \qquad (\text{constant pitch } p=b).
\end{equation}

\begin{itemize}[leftmargin=1.5em]
  \item Pitch: \(p(\vartheta)=b\).
  \item Curvature: \(\displaystyle \kappa(\vartheta)
           = \frac{b}{\sqrt{1+b^{2}}}\,e^{-b\vartheta}
           \;\xrightarrow[\vartheta\to\infty]{}\;0.\)
\end{itemize}

\begin{equation}
  \boxed{\;\omega_{\min}=0 \quad (\text{formally})\;}
\end{equation}

\emph{Interpretation:} a logarithmic spiral can extend arbitrarily far; coherence is noise-limited, not geometry-limited.

\subsubsection*{Archimedean spiral}

\begin{equation}
  r(\vartheta) = a + b\,\vartheta \qquad (\text{linear growth}).
\end{equation}

\begin{itemize}[leftmargin=1.5em]
  \item Pitch: \(\displaystyle p(\vartheta)=\frac{b}{\,a+b\vartheta\,}.\)
  \item Curvature: \(\displaystyle
    \kappa(\vartheta)
    = \frac{a + 2b\vartheta}
           {\bigl[(a+b\vartheta)^{2}+b^{2}\bigr]^{3/2}}.\)
\end{itemize}

For a finite disc of radius \(R\), the lowest sustainable rotation scales as
\begin{equation}
  \boxed{\;
    \omega_{\min}
    \;=\; C\,\frac{v\,b^{3}}{(R - b)^{2}}
    \;\propto\; R^{-2}
  \;}
\end{equation}
with \(C>0\) a shape-dependent constant.

\subsubsection*{Power-law spiral}

\begin{equation}
  r(\vartheta) = a\,\vartheta^{\,n}, \qquad n>0.
\end{equation}

\begin{itemize}[leftmargin=1.5em]
  \item Pitch: \(p(\vartheta)=\dfrac{d\ln r}{d\vartheta}=\dfrac{n}{\vartheta}\).
  \item Curvature (scaling): \(\kappa(\vartheta)\propto \vartheta^{-(n+1)}\).
\end{itemize}

Given an outer radius \(R\), the outer turn is
\begin{equation}
  \vartheta_{\max}=\Bigl(\tfrac{R}{a}\Bigr)^{\!1/n}.
\end{equation}
If \(v\) is the characteristic speed along the spiral, the low–frequency cut-off scales as
\begin{equation}
  \boxed{\;
    \omega_{\min}
    \;=\; C_{n}\,\frac{v\,a^{\,1+1/n}}{R^{\,1+1/n}}
    \;},
\end{equation}
where \(C_n>0\) is a shape-dependent constant. Note: as \(n\to 0^{+}\), \(r(\vartheta)\to a\) (a circle); at \(n=1\) one recovers the Archimedean spiral \(r=a\vartheta\). A logarithmic spiral belongs to the distinct family \(r=a\,e^{b\vartheta}\).

\subsubsection*{Composite spirals}

The composer \(\Composer\{\cdot\}\) joins different spirals at a matching radius. The composite’s cut-off is controlled by the largest \(\omega_{\min}\) among its components; a tight core (larger local \(\omega_{\min}\)) can stabilize a looser outer arm.

\subsubsection*{Practical rules of thumb}

\begin{table}[H]
  \small
  \centering
  \caption{Practical scaling of \(\omega_{\min}\).}
  \label{tab:spiral-thumb}
  \renewcommand{\arraystretch}{1.15}
  \begin{tabularx}{\linewidth}{@{}p{0.20\linewidth} p{0.28\linewidth} p{0.22\linewidth} X@{}}
    \toprule
    \textbf{Spiral family}
      & \textbf{Pitch \(p(\vartheta)\)}
      & \(\boldsymbol{\omega_{\min}}\) \textbf{scaling}
      & \textbf{Memory reach} \\
    \midrule
    Logarithmic   & constant                    & \(\to0\)                   & unlimited (noise-limited) \\
    Archimedean   & \(\propto 1/(a+\vartheta)\) & \(\propto R^{-2}\)         & finite, quadratic loss \\
    Power-law     & \(\propto 1/\vartheta\)     & \(\propto R^{-(1+1/n)}\)   & tunable via \(n\) \\
    Composite     & piece-wise                  & max of parts               & engineer-to-spec \\
    \bottomrule
  \end{tabularx}
\end{table}

\subsubsection*{Why these cut-offs matter}

\begin{itemize}[leftmargin=1.5em]
  \item \textbf{Lasers} – Choosing an Archimedean vs.\ logarithmic aperture fixes the cavity’s frequency floor.
  \item \textbf{Neural cortex} – Scroll-wave pitch sets the lowest theta rhythm that can remain phase-locked to local gamma.
  \item \textbf{Galaxies} – The power-law index of spiral arms predicts how far coherent density waves propagate before shear dominates.
\end{itemize}

\subsubsection*{Why different spiral shapes emerge}

\begin{table}[H]
  \small
  \centering
  \caption{Drivers of spiral selection.}
  \label{tab:spiral-drivers}
  \renewcommand{\arraystretch}{1.15}
  \begin{tabularx}{\linewidth}{@{}p{0.22\linewidth} p{0.33\linewidth} X@{}}
    \toprule
    \textbf{Driver}
      & \textbf{Physical meaning}
      & \textbf{Selection effect} \\
    \midrule
    Boundary conditions
      & Size/geometry of the domain (optical cavity, cortical sheet, galactic disc).
      & Finite or hard edges favour Archimedean/moderate power-law spirals; unbounded media allow logarithmic spirals. \\
    Memory-kernel dispersion
      & Strength and duration of history dependence (kernel width \(\tau\)).
      & Long, strong memory penalizes pitch variation \(\Rightarrow\) logarithmic; short, weak memory tolerates curvature change \(\Rightarrow\) power-law or composite. \\
    Energy vs.\ information budget
      & Curvature stores more “bits per turn” but costs more kinetic energy.
      & Energy-scarce, information-dense systems tighten pitch \(\Rightarrow\) nearly logarithmic; energy-rich, dissipative systems loosen pitch \(\Rightarrow\) Archimedean/power-law. \\
    \bottomrule
  \end{tabularx}
\end{table}

Formally, the spiral shape minimizes the functional
\begin{equation}
  S[p(\vartheta)]
  = \int \bigl[T_{\mathrm{kin}}(p) + V_{\mathrm{mem}}(p)\bigr]\,d\vartheta,
\end{equation}
where \(T_{\mathrm{kin}}\) rises with curvature and \(V_{\mathrm{mem}}\) penalizes rapid pitch changes. The Euler–Lagrange equation \(\delta S/\delta p=0\) yields the constant-pitch (logarithmic), linear-pitch (Archimedean), or power-law solutions depending on kernel and boundary coefficients.

%---------------------------------------------------------------
\section*{Spiral curvature, amplitude–frequency coupling, and the upper resonance limit}

The previous section fixed the \emph{lower} end of a spiral’s viable frequency window (\(\omega \ge \omega_{\min}\)). The \emph{upper} limit is set by the \RCM\ amplitude–frequency law
\begin{equation}
  A(\vartheta)\,\omega(\vartheta) \;=\; \text{const} \equiv K,
  \label{eq:amp-freq-law}
\end{equation}
which encodes energy conservation inside a resonance bundle: when the local rotation rate \(\omega\) rises, amplitude \(A\) must fall to keep the product \(K\) constant. Because curvature \(\kappa(\vartheta)\) increases the centrifugal strain on the bundle, there is a point where the required amplitude falls below the physical noise floor \(A_{\text{noise}}(\vartheta)\); coherence then collapses. That threshold defines
\begin{equation}
  \boxed{\;\omega_{\max}(\vartheta)
    \;=\; \frac{K}{A_{\text{noise}}(\vartheta)}\;}.
  \label{eq:omega-max}
\end{equation}
Combining \eqref{eq:omega-min} and \eqref{eq:omega-max} yields a finite \emph{resonance window}
\[
  \omega_{\min}(\vartheta)\;\le\;\omega(\vartheta)\;\le\;\omega_{\max}(\vartheta).
\]

Below we compute \(\omega_{\max}\) for each spiral family using a simple noise-floor model \(A_{\text{noise}}\propto\kappa^{-1/2}\) (higher curvature demands tighter confinement and therefore lower tolerable noise).

\subsubsection*{Logarithmic spiral}

\[
  r = a\,e^{b\vartheta}, \quad p=b,\quad
  \kappa(\vartheta)=\kappa_{0}\,e^{-b\vartheta},\quad
  A_{\text{noise}}\propto e^{\,b\vartheta/2}.
\]
Hence
\[
  \omega_{\max}(\vartheta)
  = \frac{K}{A_{0}}\,e^{-\,b\vartheta/2}.
\]
Since \(\omega_{\min}\propto e^{-b\vartheta}\) falls faster, a logarithmic spiral retains a non-zero bandwidth out to arbitrarily large \(\vartheta\).

\subsubsection*{Archimedean spiral}

\[
  r = a + b\vartheta,\quad
  \kappa(\vartheta)\sim\frac{a+2b\vartheta}{(a+b\vartheta)^{3}},\quad
  A_{\text{noise}}\propto(a+b\vartheta)^{3/2}.
\]
Thus
\[
  \omega_{\max}(\vartheta)
  \propto\frac{K}{(a+b\vartheta)^{3/2}}
  \;\Longrightarrow\;
  \omega_{\max}(R)\propto R^{-3/2},
\]
so the Archimedean bandwidth shrinks as \(R^{-3/2}\) and pinches out beyond a finite radius.

\subsubsection*{Power-law spiral}

\[
  r = a\,\vartheta^{n},\quad
  \kappa(\vartheta)\propto\vartheta^{-(n+1)},\quad
  A_{\text{noise}}\propto\vartheta^{(n+1)/2}.
\]
Hence
\[
  \omega_{\max}(\vartheta)\propto\vartheta^{-(n+1)/2}
  \;\Longrightarrow\;
  \omega_{\max}(R)\propto R^{-\tfrac{n+1}{2n}}.
\]
For \(n<1\) the bandwidth survives to larger radius than the Archimedean case; for \(n>1\) it shrinks faster.

\subsubsection*{Composite spirals}

For a tight–loose composite, the usable window is
\[
  [\,\omega_{\min}^{\mathrm{core}},\,\omega_{\max}^{\mathrm{arm}}\,].
\]
Designers can engineer a desired bandwidth by tuning core pitch versus outer pitch.

\subsubsection*{Bandwidth summary}

\begin{table}[H]
  \small
  \centering
  \caption{Scaling of \(\omega_{\min}\) and \(\omega_{\max}\) with radius \(R\).}
  \label{tab:bandwidth-summary}
  \renewcommand{\arraystretch}{1.15}
  \begin{tabularx}{\linewidth}{@{}p{0.20\linewidth} p{0.25\linewidth} p{0.25\linewidth} X@{}}
    \toprule
    \textbf{Spiral}      & \(\omega_{\min}(R)\)     & \(\omega_{\max}(R)\)      & \textbf{Window behaviour} \\
    \midrule
    Logarithmic         & \(0\) (formal)           & \(\propto e^{-bR/2a}\)    & Always finite, shrinks slowly \\
    Archimedean         & \(\propto R^{-2}\)       & \(\propto R^{-3/2}\)      & Pinches out at finite \(R\) \\
    Power-law           & \(\propto R^{-(1+1/n)}\) & \(\propto R^{-(n+1)/(2n)}\) & Tunable; narrows if \(n>1\) \\
    Composite           & max(core, arm)           & min(core, arm)            & Adjustable to specification \\
    \bottomrule
  \end{tabularx}
\end{table}

\subsubsection*{Physical consequences}

\begin{itemize}[leftmargin=1.5em]
  \item \textbf{Laser cavities} – Upper cut-off governs maximum stable longitudinal mode; combined with the lower cut-off it defines a “sweet spot” for gain media.
  \item \textbf{Neural oscillations} – Curvature of cortical scroll waves sets an upper gamma frequency beyond which theta–gamma phase-locking fails.
  \item \textbf{Spiral galaxies} – Outer pitch determines the highest epicyclic frequency that can remain in resonance with density waves; beyond that, arms shear apart.
\end{itemize}

\paragraph{Key takeaway.}
Every spiral class comes with a finite resonance window \([\omega_{\min},\omega_{\max}]\) fixed by its geometry, memory kernel, and noise floor. Logarithmic spirals preserve the widest bandwidth; Archimedean spirals trade robustness for a clean high-frequency ceiling; power-law and composite spirals allow intermediate compromises.

\bigskip

In the next chapter we show how these geometric limits translate into \emph{space–time coupling}, demonstrating that the same amplitude–frequency law reproduces relativistic phenomena when applied to nested spirals across scales.

The composer \(\Composer\{\cdot\}\) joins different spirals at a matching radius. The composite’s cut-off
is controlled by the largest \(\omega_{\min}\) among its components; a tight core (larger local \(\omega_{\min}\))
can stabilize a looser outer arm.


\subsubsection{Practical Rules of Thumb}

\begin{table}[H]
  \small
  \centering
  \caption{Practical Scaling of \(\omega_{\min}\)}
  \label{tab:spiral-thumb}
  \begin{tabular}{@{}p{0.18\textwidth} p{0.25\textwidth} p{0.25\textwidth} p{0.27\textwidth}@{}}
    \toprule
    \textbf{Spiral family}
      & \textbf{Pitch \(p(\theta)\)}
      & \(\boldsymbol{\omega_{\min}}\) \textbf{scaling}
      & \textbf{Memory reach} \\
    \midrule
    Logarithmic   & constant             & \(\to0\)                   & unlimited (noise-limited) \\
    Archimedean   & \(\propto1/(a+\theta)\) & \(\propto R^{-2}\)         & finite, quadratic loss \\
    Power-law     & \(\propto1/\theta\)   & \(\propto R^{-(1+1/n)}\)   & tunable via \(n\) \\
    Composite     & piece-wise            & max of parts               & engineer-to-spec \\
    \bottomrule
  \end{tabular}
\end{table}

\subsubsection{Why These Cut-offs Matter}

\begin{itemize}
  \item \textbf{Lasers} – Choosing an Archimedean vs.\ logarithmic aperture fixes the cavity’s frequency floor.
  \item \textbf{Neural cortex} – Scroll-wave pitch sets the lowest theta rhythm that can remain phase-locked to local gamma.
  \item \textbf{Galaxies} – The power-law index of spiral arms predicts how far coherent density waves propagate before shear dominates.
\end{itemize}

\subsubsection{Why Different Spiral Shapes Emerge}

\begin{table}[H]
  \small
  \centering
  \caption{Drivers of Spiral Selection}
  \label{tab:spiral-drivers}
  \begin{tabular}{@{}p{0.20\textwidth} p{0.30\textwidth} p{0.45\textwidth}@{}}
    \toprule
    \textbf{Driver}
      & \textbf{Physical meaning}
      & \textbf{Selection effect} \\
    \midrule
    Boundary conditions
      & Size/geometry of the domain (optical cavity, cortical sheet, galactic disc).
      & Finite or hard edges favour Archimedean/moderate power-law spirals; unbounded media allow logarithmic spirals. \\
    Memory-kernel dispersion
      & Strength and duration of history dependence (kernel width \(\tau\)).
      & Long, strong memory penalises pitch variation → logarithmic; short, weak memory tolerates curvature change → power-law or composite. \\
    Energy vs.\ information budget
      & Curvature stores more “bits per turn” but costs more kinetic energy.
      & Energy-scarce, information-dense systems tighten pitch → nearly logarithmic; energy-rich, dissipative systems loosen pitch → Archimedean/power-law. \\
    \bottomrule
  \end{tabular}
\end{table}

Formally, the spiral shape minimises the functional
\begin{equation}
  S[p(\theta)]
  = \int \bigl[T_{\mathrm{kin}}(p) + V_{\mathrm{mem}}(p)\bigr]\,d\theta,
\end{equation}
where \(T_{\mathrm{kin}}\) rises with curvature and \(V_{\mathrm{mem}}\) penalises rapid pitch changes.  The Euler–Lagrange equation \(\delta S/\delta p=0\) yields the constant‐pitch (logarithmic), linear‐pitch (Archimedean), or power-law solutions depending on kernel and boundary coefficients.

The next section (2.3) will show how spiral curvature also imposes an *upper* frequency bound via the amplitude–frequency law \(A\,\omega=\mathrm{const}\), completing the resonance window for each spiral family.```

\subsection{Spiral Curvature, Amplitude–Frequency Coupling, and the Upper Resonance Limit}

Sections 2.1–2.2 fixed the \emph{lower} end of a spiral’s viable frequency window (\(\omega \ge \omega_{\min}\)).  The \emph{upper} limit is set by the RCM amplitude–frequency law
\begin{equation}
  A(\theta)\,\omega(\theta) = \text{const} \equiv K,
  \tag{2.4}
\end{equation}
which encodes energy conservation inside a resonance bundle: when the local rotation rate \(\omega\) rises, amplitude \(A\) must fall to keep the product \(K\) constant.  Because curvature \(\kappa(\theta)\) increases the centrifugal strain on the bundle, there is a point where the required amplitude falls below the physical noise floor \(A_{\text{noise}}(\theta)\); coherence then collapses.  That threshold defines
\begin{equation}
  \boxed{\;\omega_{\max}(\theta)
    = \frac{K}{A_{\text{noise}}(\theta)}\;}.
  \tag{2.5}
\end{equation}
Combining (2.3) and (2.5) yields a finite \emph{resonance window}
\[
  \omega_{\min}(\theta)\;\le\;\omega\;\le\;\omega_{\max}(\theta).
\]

Below we compute \(\omega_{\max}\) for each spiral family using a simple noise‐floor model \(A_{\text{noise}}\propto\kappa^{-1/2}\) (higher curvature demands tighter confinement and therefore lower tolerable noise).

\subsubsection{Logarithmic Spiral}

\[
  r = a\,e^{b\theta}, \quad p=b,\quad
  \kappa(\theta)=\kappa_{0}\,e^{-b\theta},\quad
  A_{\text{noise}}\propto e^{\,b\theta/2}.
\]
Hence
\[
  \omega_{\max}(\theta)
  = \frac{K}{A_{0}}\,e^{-\,b\theta/2}.
\]
Since \(\omega_{\min}\propto e^{-b\theta}\) falls faster, a logarithmic spiral retains a non‐zero bandwidth out to arbitrarily large \(\theta\).

\subsubsection{Archimedean Spiral}

\[
  r = a + b\theta,\quad
  \kappa(\theta)\sim\frac{a+2b\theta}{(a+b\theta)^{3}},\quad
  A_{\text{noise}}\propto(a+b\theta)^{3/2}.
\]
Thus
\[
  \omega_{\max}(\theta)
  \propto\frac{K}{(a+b\theta)^{3/2}}
  \;\Longrightarrow\;
  \omega_{\max}(R)\propto R^{-3/2},
\]
so the Archimedean bandwidth shrinks as \(R^{-3/2}\) and pinches out beyond a finite radius.

\subsubsection{Power‐law Spiral}

\[
  r = a\,\theta^{n},\quad
  \kappa(\theta)\propto\theta^{-(n+1)},\quad
  A_{\text{noise}}\propto\theta^{(n+1)/2}.
\]
Hence
\[
  \omega_{\max}(\theta)\propto\theta^{-(n+1)/2}
  \;\Longrightarrow\;
  \omega_{\max}(R)\propto R^{-\tfrac{n+1}{2n}}.
\]
For \(n<1\) the bandwidth survives to larger radius than the Archimedean case; for \(n>1\) it shrinks faster.

\subsubsection{Composite Spirals}

For a tight–loose composite (cf.\ §1.5), the usable window is
\[
  [\,\omega_{\min}^{\mathrm{core}},\,\omega_{\max}^{\mathrm{arm}}\,].
\]
Designers can engineer a desired bandwidth by tuning core pitch versus outer pitch.

\subsubsection{Bandwidth Summary}

\begin{table}[H]
  \small
  \centering
  \caption{Scaling of \(\omega_{\min}\) and \(\omega_{\max}\) with radius \(R\)}
  \label{tab:bandwidth-summary}
  \begin{tabular}{@{}p{0.20\textwidth} p{0.25\textwidth} p{0.25\textwidth} p{0.30\textwidth}@{}}
    \toprule
    \textbf{Spiral}      & \(\omega_{\min}(R)\)     & \(\omega_{\max}(R)\)      & \textbf{Window behaviour}            \\
    \midrule
    Logarithmic         & \(0\) (formal)           & \(\propto e^{-bR/2a}\)    & Always finite, shrinks slowly       \\
    Archimedean         & \(\propto R^{-2}\)       & \(\propto R^{-3/2}\)      & Pinches out at finite \(R\)         \\
    Power‐law           & \(\propto R^{-(1+1/n)}\) & \(\propto R^{-(n+1)/(2n)}\)& Tunable; narrows if \(n>1\)         \\
    Composite           & max(core, arm)           & min(core, arm)            & Adjustable to specification         \\
    \bottomrule
  \end{tabular}
\end{table}

\subsubsection{Physical Consequences}

\begin{itemize}
  \item \textbf{Laser cavities} – Upper cut-off governs maximum stable longitudinal mode; combined with the lower cut-off (§2.2) it defines a “sweet spot” for gain media.
  \item \textbf{Neural oscillations} – Curvature of cortical scroll waves sets an upper gamma frequency beyond which theta–gamma phase‐locking fails.
  \item \textbf{Spiral galaxies} – Outer pitch determines the highest epicyclic frequency that can remain in resonance with density waves; beyond that, arms shear apart.
\end{itemize}

\paragraph{Key takeaway.}  
Every spiral class comes with a finite resonance window \([\omega_{\min},\omega_{\max}]\) fixed by its geometry, memory kernel, and noise floor.  Logarithmic spirals preserve the widest bandwidth; Archimedean spirals trade robustness for a clean high‐frequency ceiling; power‐law and composite spirals allow intermediate compromises.



In Chapter 3 we will show how these geometric limits translate into \emph{space–time coupling}, demonstrating that the same amplitude–frequency law reproduces relativistic phenomena when applied to nested spirals across scales.  


